\chapter{Results and Discussion}
\label{ch:results}

This chapter reports measured results from a real experimental run of
the implemented system (Chapter~\ref{ch:impl}). Every number is
traceable to the recorded run; no value is a placeholder or an
estimate. Where a research question requires the full verifier of
Chapter~\ref{ch:method}, which is not yet built, the question is
answered as far as the implemented main--branch verifier (presence and
value grounding) allows and the remainder is marked as future work.

\section{Experimental Setup}

The evaluation is a controlled \textbf{270--cell sweep}. The unit of
analysis is a cell --- a tuple
\[
(\text{dataset}, \text{condition}, \text{LLM backend}, \text{seed}),
\]
and the design is the full cross--product of \textbf{6 datasets
$\times$ 3 conditions $\times$ 3 backends $\times$ 5 seeds}. Each cell
produces one \texttt{run\_results} row.

\begin{itemize}
    \item \textbf{Conditions.} B0 (naive), B1 (schema), B2
    (meta--feature--guided), as implemented in
    Chapter~\ref{ch:impl}.
    \item \textbf{Backends.} \texttt{gpt-oss:120b-cloud},
    \texttt{gemma4:31b-cloud} and \texttt{nemotron-3-nano:30b-cloud},
    served through Ollama.
    \item \textbf{Seeds.} 3001--3005 (five), fixing the train/test
    split and generation seed so comparisons across conditions are
    paired on the same cell.
    \item \textbf{Metric.} All six datasets are binary
    classification; the score is \texttt{balanced\_accuracy} on a
    held--out test split (higher is better), chosen so that a
    majority--class predictor cannot score well.
    \item \textbf{Max iterations.} 3; a cell is a failure if no
    runnable, verified pipeline is produced within three attempts.
\end{itemize}

The run completed in 5\,h\,17\,m. A leakage guard (any
\texttt{balanced\_accuracy} $\geq 0.995$ on a classification task)
fired on \textbf{zero} cells. All artefacts (generated code, traces,
verification reports) are persisted, so the tables below can be
regenerated from storage without re--calling any LLM.

\begin{quote}
\textbf{Note on earlier data.} An initial mid--term snapshot
(2026--07) contained only six runs under a single condition and is
\emph{not} used here; it has been fully superseded by the 270--cell
sweep. A first sweep attempt was also discarded after a mid--run
reboot corrupted it; the reported run is the clean redo.
\end{quote}

\subsection{Datasets}

\begin{table}[H]
\centering
\caption{The six--dataset panel (all binary classification).}
\label{tab:datasets}
\small
\renewcommand{\arraystretch}{1.2}
\begin{tabular}{l r r l}
\toprule
\textbf{Dataset} & \textbf{Rows} & \textbf{Cols} & \textbf{Source} \\
\midrule
wdbc        & 569  & 31 & OpenML 1510 \\
breast-w    & 699  & 10 & OpenML 15 \\
diabetes    & 768  & 9  & OpenML 37 \\
Titanic     & 891  & 11 & curated \\
spambase    & 4601 & 58 & OpenML 44 \\
Telco churn & 7043 & 21 & curated (IBM) \\
\bottomrule
\end{tabular}
\end{table}

\section{Pipeline-Level Results}

Table~\ref{tab:percond} gives the headline per--condition result
collapsed over all datasets, backends and seeds (90 cells each).

\begin{table}[H]
\centering
\caption{Per--condition results over all 90 cells each. ``Mean score''
and ``Mean iters'' are over successful runs only.}
\label{tab:percond}
\renewcommand{\arraystretch}{1.3}
\begin{tabular}{l r r r r}
\toprule
\textbf{Condition} & \textbf{Success} & \textbf{Success \%} &
\textbf{Mean score} & \textbf{Mean iters} \\
\midrule
B0 (naive)        & 90/90 & 100.0\,\% & 0.8432 & 1.01 \\
B1 (schema)       & 90/90 & 100.0\,\% & 0.8468 & 1.06 \\
B2 (meta-feature) & 56/90 & \textbf{62.2\,\%} & 0.8549$^\ast$ & 1.41 \\
\bottomrule
\end{tabular}
\end{table}

\noindent$^\ast$ B2's marginally higher mean is \emph{survivorship
bias}: only the 56 cells it completes contribute, while its 34
failures (often the harder cells) are excluded. The paired test in
Section~\ref{sec:rq4} is the valid comparison.

The breakdown by backend (Table~\ref{tab:permodel}) shows the effect
is uniform and that B2 is far more expensive. B0 and B1 are perfectly
reliable across all three backends; B2 drops every backend and is
catastrophic for the smallest (nemotron--nano, 26.7\,\%), whose
run--away completion length (26k tokens) trips upstream limits.

\begin{table}[H]
\centering
\caption{Per--backend $\times$ condition: success, iterations
(averaged over all 30 cells, including failed attempts) and token
cost (averages).}
\label{tab:permodel}
\scriptsize
\renewcommand{\arraystretch}{1.25}
\begin{tabular}{l l r r r r r}
\toprule
\textbf{Backend} & \textbf{Cond} & \textbf{ok/30} & \textbf{succ\%} &
\textbf{iters} & \textbf{prompt tok} & \textbf{compl.\ tok} \\
\midrule
gemma4:31b-cloud          & B0 & 30 & 100.0 & 1.00 & 348  & 627 \\
gemma4:31b-cloud          & B1 & 30 & 100.0 & 1.00 & 1182 & 611 \\
gemma4:31b-cloud          & B2 & 24 & 80.0  & 1.87 & 9280 & 2506 \\
gpt-oss:120b-cloud        & B0 & 30 & 100.0 & 1.00 & 369  & 1338 \\
gpt-oss:120b-cloud        & B1 & 30 & 100.0 & 1.03 & 1170 & 1346 \\
gpt-oss:120b-cloud        & B2 & 24 & 80.0  & 1.70 & 6449 & 5330 \\
nemotron-3-nano:30b-cloud & B0 & 30 & 100.0 & 1.03 & 348  & 2112 \\
nemotron-3-nano:30b-cloud & B1 & 30 & 100.0 & 1.13 & 1433 & 2660 \\
nemotron-3-nano:30b-cloud & B2 & 8  & \textbf{26.7} & 1.33 & 7373 & \textbf{26244} \\
\bottomrule
\end{tabular}
\end{table}

\begin{figure}[H]
    \centering
    \miroplaceholder{figures/project_figures/fig_success_by_condition.png}{%
    Grouped bar chart of success rate (\%) by condition, one group per
    backend; B0 and B1 at 100\,\% and B2 at 80/80/27\,\%. A second
    panel can overlay mean prompt tokens to show B2's cost.}
    \caption{Success rate by condition and backend.}
    \label{fig:success}
\end{figure}

Table~\ref{tab:perdataset} reports mean test score per dataset and
condition. On every dataset the three conditions land within $\pm
0.02$ of each other; meta--features do not shift pipeline quality.

\begin{table}[H]
\centering
\caption{Mean \texttt{balanced\_accuracy} per dataset $\times$
condition (successful runs; ``ok'' = successes of 15).}
\label{tab:perdataset}
\scriptsize
\renewcommand{\arraystretch}{1.25}
\begin{tabular}{l r | r r | r r | r r}
\toprule
& & \multicolumn{2}{c|}{\textbf{B0}} & \multicolumn{2}{c|}{\textbf{B1}}
& \multicolumn{2}{c}{\textbf{B2}} \\
\textbf{Dataset} & \textbf{rows} & ok & score & ok & score & ok & score \\
\midrule
wdbc        & 569  & 15 & 0.9618 & 15 & 0.9673 & 12 & 0.9707 \\
breast-w    & 699  & 15 & 0.9524 & 15 & 0.9492 & 10 & 0.9554 \\
diabetes    & 768  & 15 & 0.7215 & 15 & 0.7172 & 11 & 0.7126 \\
Titanic     & 891  & 15 & 0.7874 & 15 & 0.7898 & 6  & 0.7670 \\
spambase    & 4601 & 15 & 0.9352 & 15 & 0.9313 & 10 & 0.9350 \\
Telco churn & 7043 & 15 & 0.7009 & 15 & 0.7259 & 7  & 0.6974 \\
\bottomrule
\end{tabular}
\end{table}

\section{Evaluation of the Research Questions}

\subsection{RQ1 --- Representation}

The implemented B2 pipeline lifts each plan into a structured
reasoning trace (a list of typed \texttt{Decision} claims) that is
\emph{mechanically} comparable with the generated code. All 90 B2
cells produced a parseable trace that the verifier consumed without a
human in the loop. This is preliminary support for RQ1 \emph{at the
level of the implemented representation} (presence + value grounding);
the richer graph--based PIR of Chapter~\ref{ch:method} remains future
work, so full coverage of ordering and dataflow claims is not yet
demonstrated.

\subsection{RQ2 --- Discrimination (and H1, H2)}
\label{sec:rq2}

This is the project's central claim and it is supported on real data.
The verifier flagged \textbf{seven} B2 runs as
\texttt{reasoning\_unfaithful}. \emph{All seven executed successfully
and produced good scores} (0.717, 0.721, 0.722, 0.925, 0.942, 0.944,
0.969) --- yet their stated reasoning did not match their code. An
execution--only or performance--only check would have \textbf{passed
all seven}; the mechanical check rejected them. This is a direct,
if small, confirmation of \textbf{H1} (execution success does not
imply alignment) and preliminary support for \textbf{H2} (the
mechanical check discriminates where execution cannot).

\subsection{RQ3 --- Categories}

Only B2 produced failures; B0 and B1 were 100\,\% runnable. The
observed failure categories are shown in Table~\ref{tab:errors}. The
implemented verifier emits a single misalignment bucket,
\texttt{reasoning\_unfaithful}; the finer taxonomy of
Chapter~\ref{ch:method} (missing operation, wrong parameter, wrong
order, \ldots) is designed but not yet emitted, so RQ3 is answered
only partially --- misalignment \emph{does} occur in practice (7
cases), but its sub--categorisation is future work.

\begin{table}[H]
\centering
\caption{Failure categories (B2 only; B0 and B1 had none).}
\label{tab:errors}
\small
\renewcommand{\arraystretch}{1.2}
\begin{tabular}{l l r}
\toprule
\textbf{Backend} & \textbf{Category} & \textbf{n} \\
\midrule
gemma4:31b-cloud          & reasoning\_unfaithful & 5 \\
gemma4:31b-cloud          & syntax\_error         & 1 \\
gpt-oss:120b-cloud        & runtime\_other        & 4 \\
gpt-oss:120b-cloud        & api\_hallucination    & 1 \\
gpt-oss:120b-cloud        & reasoning\_unfaithful & 1 \\
nemotron-3-nano:30b-cloud & infrastructure        & 19 \\
nemotron-3-nano:30b-cloud & import\_error         & 2 \\
nemotron-3-nano:30b-cloud & reasoning\_unfaithful & 1 \\
\bottomrule
\end{tabular}
\end{table}

\begin{figure}[H]
    \centering
    \miroplaceholder{figures/project_figures/fig_error_taxonomy.png}{%
    Stacked bar / treemap of the failure categories in
    Table~\ref{tab:errors}, grouped by backend, highlighting the seven
    \texttt{reasoning\_unfaithful} detections and nemotron--nano's 19
    infrastructure failures.}
    \caption{Distribution of B2 failure categories.}
    \label{fig:errtax}
\end{figure}

\subsection{RQ4 --- Variation and Predictive Score}
\label{sec:rq4}

\textbf{Variation.} Faithfulness and reliability vary clearly by
backend: all seven unfaithful flags concentrate as gemma~5,
gpt-oss~1, nemotron~1, and B2 success ranges from 80\,\% (the two
larger models) to 26.7\,\% (nemotron--nano). This is preliminary
support for RQ4's claim that behaviour varies with the LLM; variation
by pipeline complexity (H4) is \emph{not} tested here, as the panel
was not stratified by operation count.

\textbf{Predictive score.} A paired Wilcoxon signed--rank test on
\texttt{test\_score}, matched on each \((\text{dataset}, \text{seed},
\text{backend})\) cell where both conditions succeeded
(Table~\ref{tab:wilcoxon}), shows \textbf{no significant difference}
for any pair (Bonferroni--adjusted $\alpha = 0.0167$). Meta--feature
prompting does not improve predictive quality over naive or schema
prompting.

\begin{table}[H]
\centering
\caption{Paired Wilcoxon signed--rank on test score (both conditions
succeeded on the matched cell). Median $\Delta$ is the second
condition minus the first.}
\label{tab:wilcoxon}
\small
\renewcommand{\arraystretch}{1.25}
\begin{tabular}{l r r r r l}
\toprule
\textbf{Pair} & \textbf{n} & \textbf{median $\Delta$} & \textbf{z} &
\textbf{p} & \textbf{verdict} \\
\midrule
B0 vs B1 & 64 & $+0.0044$ & $-1.669$ & 0.095 & n.s. \\
B0 vs B2 & 54 & $-0.0022$ & $-0.637$ & 0.524 & n.s. \\
B1 vs B2 & 51 & $+0.0043$ & $-0.356$ & 0.722 & n.s. \\
\bottomrule
\end{tabular}
\end{table}

\begin{figure}[H]
    \centering
    \miroplaceholder{figures/project_figures/fig_score_comparison.png}{%
    Box/strip plot of per--cell test scores grouped by condition
    (optionally faceted by dataset), making visible that the three
    conditions overlap and no pair separates --- the Wilcoxon result
    of Table~\ref{tab:wilcoxon}.}
    \caption{Predictive score by condition.}
    \label{fig:scores}
\end{figure}

\subsection{RQ5 --- Localisation}
\label{sec:rq5}

The implemented verifier localises a failure to the \emph{specific
decision(s)} that did not verify (it records, per decision, whether
the action was present and whether the value matched, with a
human--readable reason). It does not yet localise to a \emph{source
span} in the code, which requires the source--located PIR of
Chapter~\ref{ch:method}. RQ5 is therefore supported at claim level and
pending at span level.

\section{Hypothesis Summary}

\begin{table}[H]
\centering
\caption{Status of the four hypotheses against the 270--cell run.}
\label{tab:hyp}
\small
\renewcommand{\arraystretch}{1.3}
\begin{tabular}{l p{8.5cm} l}
\toprule
\textbf{H} & \textbf{Statement} & \textbf{Status} \\
\midrule
H1 & Execution success does not imply alignment & \textbf{Supported}
(7 executing, scoring runs flagged unfaithful) \\
H2 & Mechanical alignment discriminates better than execution success
& Preliminary support \\
H3 & Mismatches form stable AutoML--specific categories & Pending
(full taxonomy not yet emitted) \\
H4 & Mismatch probability grows with pipeline complexity & Pending (no
complexity stratification) \\
\bottomrule
\end{tabular}
\end{table}

\section{Comparison with Baselines (Planned)}

Chapter~\ref{ch:method} specifies a comparison of the mechanical
verifier against four baselines --- execution--only,
performance--only, LLM--judge and keyword/API --- on a
seeded--mutation corpus with human--annotated targets. That corpus
and its annotations have not yet been built, so the quantitative
comparison (false--alignment rate, category macro--F1, localisation
IoU) is \textbf{future work}; Table~\ref{tab:results} records its
pre--committed structure.

\begin{table}[H]
\centering
\caption{Structure of the planned baseline comparison. Entries are
left blank pending the seeded--mutation corpus.}
\label{tab:results}
\renewcommand{\arraystretch}{1.2}
\begin{tabular}{lcccc}
\toprule
\textbf{Method} & \textbf{Coverage} & \textbf{S\textsubscript{align}}
& \textbf{CMR $\downarrow$} & \textbf{Loc.\ macro--F1} \\
\midrule
Execution--only & -- & -- & -- & -- \\
Performance--only & -- & -- & -- & -- \\
LLM judge & -- & -- & -- & -- \\
Keyword / API match & -- & -- & -- & -- \\
Mechanical verifier (proposed) & -- & -- & -- & -- \\
\bottomrule
\end{tabular}
\end{table}

One claim must be stated carefully. It is \emph{definitionally} true
that an execution-- or performance--only baseline, which never reads
the plan, cannot flag a reasoning--code misalignment that is
compatible with running code and a good score --- the seven runs of
Section~\ref{sec:rq2} are concrete instances. But a precise
\emph{false--alignment rate} for such a baseline is a quantity to be
\emph{measured} on the annotated corpus, not asserted as 100\,\%; the
earlier draft's blanket ``100\,\% false alignment'' claim has been
removed as unsupported.

\section{Discussion}

\textbf{A useful negative result.} Across three backends and six
datasets, injecting meta--features and decision rules into the prompt
(B2) did not improve predictive score (Table~\ref{tab:wilcoxon}) and
markedly reduced reliability (100\,\% $\rightarrow$ 62\,\%) at
5--25$\times$ the prompt--token cost. For this family of models and
tabular tasks, the schema signal (B1) already captures what helps; the
extra meta--feature payload mostly adds failure surface. This is a
clean, reportable finding, not a disappointment --- it tells a
practitioner they can stop paying the meta--feature cost.

\textbf{The verifier earns its place.} The same B2 machinery that
hurts reliability is what makes faithfulness \emph{checkable}. The
seven unfaithful detections are exactly the cases execution-- and
performance--based evaluation cannot see, which is the thesis of this
project demonstrated end to end on real LLM output.

\textbf{A caveat on the trace.} Because the implemented B2 trace is
printed by the generated code at runtime, it is a post--hoc
self--report rather than an independently captured plan. The verifier
can therefore catch a trace that contradicts its own code, but a model
that writes mismatched code \emph{and} a matching (equally wrong)
trace would pass. Independent pre--code reasoning capture
(\texttt{B2R}) is the key next step.

\section{Threats to Validity}

\textbf{External validity.} Three Ollama backends and six binary
classification datasets are a small sample; results may not
generalise to other model families, to regression or multiclass
tasks, or to high--cardinality/time--series data.

\textbf{Construct validity.} B0/B1/B2 are three specific prompts, not
the best possible prompt in each class; a different B2 wording could
behave differently. The faithfulness verifier checks presence and
value grounding only --- it does not yet check parameters, ordering or
columns, so ``faithful'' here is a necessary, not sufficient,
condition for full alignment.

\textbf{Internal validity.} The implemented verifier treats an
\emph{uncheckable} value as not--contradicted, so a decision citing no
numeric meta--feature cannot currently be refuted; the proposed
three--valued verdict would mark it \textsc{unverifiable}. Survivorship
bias affects any score comparison that conditions on B2 success, which
is why the paired test restricts to matched, both--succeeded cells.

\textbf{Conclusion validity.} Five seeds and a single seed family
(3001--3005) bound the statistical power; the non--significant score
differences should be read as ``no evidence of an effect'' at this
sample size, not as proof of exact equality.

\textbf{Reproducibility.} Model identifiers, seeds, generated code and
JSON traces are all persisted, and the analysis layer reads only
stored rows, so every table here is reproducible from the archived
run.

\section{Summary}

On a real 270--cell sweep, the naive and schema prompt conditions were
perfectly reliable while the meta--feature condition was markedly less
reliable and far more expensive without improving predictive score
(paired Wilcoxon, all $p>0.09$). The implemented mechanical verifier
flagged seven pipelines that executed and scored well but whose
reasoning did not match their code --- supporting H1 and giving
preliminary support to H2, the project's core hypotheses. The
representation (RQ1), discrimination (RQ2) and claim--level
localisation (RQ5) questions are answered at the level of the
implemented main--branch verifier (presence + value grounding); the
full taxonomy (RQ3/H3), complexity
variation (H4) and the quantitative baseline comparison remain future
work pending the richer verifier and a seeded--mutation corpus.
